% A coefficient estimate and its band consequence; not the general conjecture.
\section{A midpoint estimate for the penultimate unresolved difference}
For two nonrectangular increasing arms put $P=F_XF_Y$, $D=\deg P$,
$q_j=P_{D-j}$, and $r=r_X+r_Y$. Write the arms as
$1^\alpha U,1^\beta V$, with positive initial strings and all remaining
parts at least two. Let $h$ be the number of factors of length at least
three in the main reverse box
$[\alpha+1][\beta+1]\prod_{t\in U,V}[t]$.
Deleted-prefix terms begin at degree at least three, so
\[
 q_1=r,\qquad q_2=\binom r2+h,
 \qquad 2q_2-q_1=r(r-2)+2h\le r^2.
\]

\begin{lemma}[Central coefficients are capacity-monotone]
Increasing any initial-string length or any non-one part cannot decrease
$P_{\lfloor D/2\rfloor}$. The non-one lists need not be sorted here.
\end{lemma}
\begin{proof}
The forward append recurrence shows that increasing any parameter
increases the polynomial coefficientwise. It suffices to increase a
parameter by one. If the middle index stays fixed, the claim is immediate.
If it increases by one, the proved ordinary-product growth through
$\lfloor D_{\rm new}/2\rfloor+1$ shows that the new middle coefficient
is at least its preceding coefficient, which in turn is at least the old
middle coefficient. Both factors remain nonrectangular throughout.
\end{proof}

\begin{theorem}[An exact Fibonacci lower envelope for the middle height]
For any two nonrectangular arms in the formal no-return-to-one class,
\[
           P_{\lfloor D/2\rfloor}\ge \operatorname{Fib}(r+2)-2.
\]
Equality is attained by the pair $1,2$ and $1,2^{r-3}$, for every $r\ge4$.
\end{theorem}
\begin{proof}
Reduce capacities to the minimal binary arms by the preceding lemma.
Write $A_s=F_{1,2^s}$ and $T=z+z^2$. For $1\le s\le t$,
the recurrence $A_j=1+TA_{j-1}$ gives
\[
 A_sA_t-A_{s-1}A_{t+1}=A_t-A_{s-1}\ge0
\]
coefficientwise. Repeatedly apply this identity until the smaller
non-one count is one. At fixed $s+t=r-2$, the resulting product is
$A_1A_{r-3}$ and its middle rank is $r-1$.
The arm $A_{r-3}$ has length $r-2$: its coefficients through that
length are $\operatorname{Fib}(n+1)$, while its next coefficient is
$\operatorname{Fib}(r)-2$. Since $A_1=1+z+2z^2+z^3$, extraction gives
\[
 [z^{r-1}]A_1A_{r-3}
 =\operatorname{Fib}(r)-2+\operatorname{Fib}(r-1)
       +2\operatorname{Fib}(r-2)+\operatorname{Fib}(r-3)
 =\operatorname{Fib}(r+2)-2.
\]
The formulas include $r=4$ directly. This also proves sharpness.
\end{proof}

\begin{corollary}[A whole interior interval]
For every $r-1\le n\le D-r+1$,
\[
                  P_n\ge\operatorname{Fib}(r+2)-2.
\]
\end{corollary}
\begin{proof}
Both forward coefficients and reverse coefficients are nondecreasing
under arm-capacity increases, by the proved capacity comparison.
Consequently $P_{r-1}$ and $P_{D-r+1}$ are at least the middle
coefficient of the corresponding minimal binary product, hence at
least the displayed Fibonacci number minus two. Unimodality of the
ordinary product gives the same bound between the two endpoints.
\end{proof}

\begin{corollary}[A direct terminal-difference certificate]
Let $j\ge0$, put $e=D-M$, and define
\[
 B(r,j)=2\binom{r+j}{j+1}-\binom rj.
\]
Suppose
\[
 e\ge r+j,\qquad
 \operatorname{Fib}(r+2)-2\ge B(r,j),\qquad
 \binom r{j+2}\ge B(r,j).
\]
At any such rank $D-j>\max(a,b)+1$, one has
$\Delta H(D-j)\le0$.
\end{corollary}
\begin{proof}
The established top-coefficient bounds imply
$2q_{j+1}-q_j\le B(r,j)$ and $q_{j+2}\ge\binom r{j+2}$.
The preceding proof gives $P_{r-1}\ge\operatorname{Fib}(r+2)-2$.
Also $r-1\le e-j-1\le D-j-2$, where the last inequality uses
$M\ge1$. Unimodality therefore gives
$P_{e-j-1}\ge B(r,j)$ from the endpoint bounds at ranks
$r-1$ and $D-j-2$. Substitute into the exact central difference
$\Delta H(D-j)=2q_{j+1}-q_j-P_{e-j-1}$.
This certificate compares coefficients directly and makes no
replacement of them by cumulative sums.
\end{proof}

\begin{lemma}[A large-cover central bound]
If $r\ge9$, then $P_{\lfloor D/2\rfloor}\ge r^2+1$.
\end{lemma}
\begin{proof}
By capacity monotonicity, reduce both arms to
$1,2^{r_X-1}$ and $1,2^{r_Y-1}$, keeping their cover counts fixed.
The new total degree is $2r-2$. For an individual minimal arm,
\[
 q_{r_X}=(1+z)^{r_X}
       +\sum_{i=1}^{r_X-1}z^{2i+1}(1+z)^{r_X-i-1}.
\]
Extract degree $r-1$ from the product. Keep the main term, the first
deleted-prefix term from each arm, their product, and the second and
third deleted-prefix terms from one arm having at least four covers.
Such an arm exists since $r\ge9$. All omitted terms are nonnegative.
The resulting lower bound is
\begin{align*}
 r+2\binom{r-2}{2}+\binom{r-3}{3}
       +\binom{r-4}{3}+\binom{r-4}{4}
 &=r+(r-2)(r-3)+\binom{r-2}{4}\\
 &\ge r^2+1.
\end{align*}
For the last step, $\binom{r-2}{4}\ge4r-5$ holds at $r=9$ and
persists by Pascal's identity, since its next increment is
$\binom{r-2}{3}\ge35>4$.
\end{proof}

\begin{lemma}[The finite small-cover inputs]
The following exact inequalities hold.
\begin{enumerate}
\item For $4\le r\le7$ and $D=23$, one has
      $P_{11}\ge r(r-2)+2h$.
\item For $r=8$, $\alpha,\beta\in\{1,2\}$ and all parts of
      $U,V$ in $\{2,3\}$, one has
      $P_{\lfloor D/2\rfloor}\ge r(r-2)+2h+1$.
\end{enumerate}
The arms are increasing in these two statements.
\end{lemma}
\begin{proof}
These are finite, explicitly specified inputs. Exchange the arms to
assume $2\le r_X\le r_Y$. For the first assertion enumerate
$a=3,\ldots,20$, $\alpha=1,\ldots,a$, and the nondecreasing lists
$U$ of $r_X-1$ parts at least two summing to $a-\alpha$; enumerate
the other arm at weight $23-a$. There are respectively
$1140,2160,5031,3560$ profiles for $r=4,5,6,7$.
For the second assertion there are $172$ profiles. The exact minimum
margins are respectively $0,3,14,43$ and $3$, with the extra unit
already subtracted in the last case.

\texttt{MidpointTerminalFinite.lean} implements these enumerators and
the forward recurrence, and states their exhaustive inequalities using
kernel reduction. The accompanying Python audit independently enumerates
the same finite domains. These finite inputs do not replace the
unbounded reductions in the next proof.
\end{proof}

\begin{theorem}[Midpoint coefficient pays the top two ranks]
If $D\ge23$, then
\[
                  P_{\lceil D/2\rceil+1}\ge2q_2-q_1.
\]
\end{theorem}
\begin{proof}
First compare the indicated coefficient with $P_{\lfloor D/2\rfloor}$.
If $D$ is even, ordinary-product growth gives the comparison without loss.
If $D$ is odd, every product of arm height layers other than the two
lowest layers has endpoint sum at least $D+1$, so its reflection gives
the comparison without loss. The lowest-layer rectangle loses at most
one, and loses that unit only when $|a-b|=1$.

For $r\ge9$, the large-cover bound pays this possible unit and the
target $2q_2-q_1\le r^2$. For $r=8$, clip the initial lengths to two
and non-one parts to three. This preserves $r,h$ and cannot increase
the central coefficient. The second finite input again pays the unit.

Suppose $4\le r\le7$. The fully clipped weight is at most
$3r-2\le19$. Therefore, if $D\ge23$, parameters greater than their
clipping levels can be decreased one at a time until the weight is
exactly $23$, preserving $r,h$. The non-one parts can be kept sorted
by decreasing the first part in a block of equal largest reducible
parts. The first finite input and capacity monotonicity prove the target
at $P_{\lfloor D/2\rfloor}$.

It remains only to pay the possible unit when $|a-b|=1$.
Put $a=\min(a,b)$, so $a\ge11$. Each nonrectangular arm has lowest
layer $[0,a]$ and second layer $[2,a-1]$, with the corresponding
intervals for the other arm. Their four products contribute
$4a-5$ at the middle rank $a$. If $r=4$, this exceeds $r^2+1$.
If $5\le r\le7$, one arm has at least three top covers; its third
layer contains $[3,a-2]$ (or $[3,b-2]$). Multiplying that interval
by the first two layers of the other arm adds at least $2a-8$.
Thus the middle coefficient is at least $6a-13\ge53>r^2+1$.
The possible loss of one is paid in all cases.
\end{proof}

\begin{theorem}[Three times the peak plus seven]
Every one-descent composition with descent peak $M$ and weight
$N\le3M+7$ has a unimodal rank polynomial. Both relative polynomials
of strict cores in this band are unimodal as well.
\end{theorem}
\begin{proof}
Use a least relative counterexample in the deletion- and exchange-closed
strict-core family. The proved $3M+6$ band and the peak-seven theorem
leave only $M\ge8$, balanced nonrectangular arms, and
$D=2M+7$, $e=D-M=M+7$.
The improved head bound gives growth through
$n_0=D-\lfloor e/2\rfloor+1$.

For $n_0+1\le n\le D-3$, write $j=D-n$. Then
$j\ge3$, $j+1<\min(a,b)$ and
$2j+1\le e-3=M+4=\lceil D/2\rceil$.
The deficit-layer theorem gives $\mathcal T_j(q)\ge0$.
Also $e-j\le M+4=\lfloor D/2\rfloor+1$, where $P$ is still
nondecreasing. Hence
\[
 \Delta H(n)-\Delta H(n+1)
 =\mathcal T_j(q)+P_{e-j}-P_{e-j-1}\ge0.
\]
This makes the undecided differences through rank $D-2$ antitone.

Since $D\ge23$ and $e-2=\lceil D/2\rceil+1$, the midpoint theorem gives
\[
 \Delta H(D-1)=2q_2-q_1-P_{e-2}\le0.
\]
The existing terminal lemma gives $\Delta H(D)\le0$, and all differences
from $D+1$ onward are nonpositive. The initial nonnegative segment,
antitone middle, and nonpositive final segment preclude a rebound.
Restore repeated maxima and the increasing final tail as before.
\end{proof}

Consequently a counterexample with $N-2M\le E$ must satisfy
$M\le E-8$, $N\le3E-16$. All one-descent compositions with
$N\le2M+15$, and all such compositions of total weight at most
$31$, are unimodal. The unrestricted conjecture is still open here.
